\section{KESENJANGAN PENELITIAN (\textit{RESEARCH GAP})}
\label{sec:research-gap}

% Identifikasi celah atau kesenjangan yang belum terselesaikan oleh penelitian terdahulu.
Berdasarkan analisis kritis terhadap literatur yang ditinjau, ditemukan beberapa kesenjangan penelitian (\textit{research gap}) utama:

\subsection{Kesenjangan Metodologis dan Integrasi}
[Uraikan kesenjangan metodologis, keterbatasan teknis, atau kekurangan integrasi dari penelitian terdahulu...]

\subsection{Kesenjangan Evaluasi dan Pengujian}
[Uraikan kekurangan atau celah dalam hal evaluasi empiris, dataset, atau pengujian performa...]
